\cotefile{77}
\scanpage{34}
L'ens des sections de $2F$ est \deleted{\unsure{vide ou}} \deleted{\unsure{(qui exprime|(qui est}} 
\deleted{torseur sous $\Gamma(S,L\unsure{)}$} non vide \unsure{sss|ssi} la classe $2\,\mathrm{cl}(F)=\mathrm{cl}(2F)$
dans $H^1(S,L)$ est $=0$. Alors cet ens est
un torseur sous $H^0(S,L)$ (qui exprime
la première indétermination). La section
$\underline b$ étant choisie, l'ens des $\sigma$ au dessus
de $\underline b$ est non vide sss la classe
\[
 c(\underline b)=\mathrm{cl}(\Sigma(\underline b))\in H^1(S,L^{\otimes2})
\]
est $=0$; et dans cet cas est un torseur
sous $H^{\unsure{1|0}}(S,L^{\otimes2})$.

On peut considérer que $\Sigma^2(F)$, torseur
affine sous $\check B\otimes L$, est donné par une
structure d'extension
\[
 \unsure{\cdot|0}\unsure{-|\to}\check B\otimes L\longrightarrow V^{(3)}\longrightarrow\mathcal O_S\longrightarrow0 \tag*{(45)}
\]
(où on considère $\check B$ comme $V^{(2)}$).
Or $\mathrm{Sym}^2(\check B)=\mathrm{Sym}^2(V^{(1)})$ est une
extension de $\mathrm{Sym}^2(\mathcal O_S)=\mathcal O_S$ par $\mathcal O_S\otimes V^{\unsure{(2)|(1)}}=V^{(1)}$.
En fait, la proposition précédente se
précise en
\[
\begin{aligned}
 \Sigma^2(F)={}&\text{fibré affine associé à}\\
 &\text{l'extension }V^{(3)}=\mathrm{Sym}^2(\check B)\\
 &\text{de }\mathcal O_S\text{ par }V^{\unsure{(2)|(2)'|(1)}}=\check B.
\end{aligned}
\tag*{(46)}
\]
% archive sheet number: 33

\scanpage{35}
(\uline{5 bis}) \uline{Détermination de} $\Sigma^n(F)$ (2\textsuperscript{ème} mouture \ldots)

Considérons l'extension vectorielle définissant $F$ \inserted{un fibré affine de rg 1}
\[
0\longrightarrow L\xrightarrow{i_1}\deleted{\unsure{B}}\overset{\deleted{\unsure{6|e|\leq}}}{V}\xrightarrow{\pi_1}\mathcal O_S\longrightarrow0 \qquad (V=\check B) \tag*{(47)}
\]
\[
 F=\pi_1^{-1}(1)\subset\check{\mathbb V}(V) \tag*{(48)}
\]
on en déduit, pour tout $n\geq1$, une structure
d'extension
\[
0\longrightarrow L\otimes\underline{\mathrm{Sym}}^{n-1}(\deleted{\unsure{B}}V)\xrightarrow{i_n}\underline{\mathrm{Sym}}^n(\deleted{\unsure{\check B}}V)\xrightarrow{\pi_n}\mathcal O_S\longrightarrow0 \tag*{(49)}
\]
qui \inserted{(en \unsure{itérant|théorie})} met sur $\mathrm{Sym}^n(V)$ une filtration
\deleted{\unsure{dont les}} décroissante par les
\[
\underbrace{L^i\otimes\underline{\mathrm{Sym}}^{n-i}(V)}_{\mathrm{Fil}^i\,\mathrm{Sym}^n(V)}\subset\underline{\mathrm{Sym}}^n(V)\quad(0\leq i\leq n) \tag*{(50)}
\]

\deleted{Considérons}
\[
\deleted{\Sigma^n\deleted{\unsure{(F)|F}}=\pi_n^{-1}(\deleted{\unsure{\check{\mathbb V}(\mathcal O_S)}}1)\subset\check{\mathbb V}(\underline{\mathrm{Sym}}^n(V))} \tag*{\deleted{(51)}}
\]
\deleted{le fibré affine associé à l'extension (50).}
\deleted{On va définir un isomorphisme canonique}
\[
\deleted{\underbrace{\Sigma^n(F)}_{\substack{n.\text{ième puiss.}\\\text{symétrique}}}\xrightarrow[\sim]{\varphi_n}\Sigma^n} \tag*{\deleted{(52)}}
\]

\deleted{\unsure{1)|6)} Description ``fonctorielle''. On définira le
morphisme $\varphi_n$ par passage au quotient du}
\[
\deleted{\left\{\begin{aligned}
 \psi_n &: F^n\longrightarrow\Sigma^n\\
 \psi_n(x_1,\ldots,x_n)&=x_1
\end{aligned}\right.} \tag*{\deleted{(53)}}
\]
% archive sheet number: 34

\scanpage{36}
Considérons aussi la suite exacte transposée de (47)
\[
0\longrightarrow\mathcal O_S\xrightarrow{i'_1}\underset{\substack{\Vert\\ B}}{\check V}\xrightarrow{\pi'_1}\check L\longrightarrow0 \tag*{(51)}
\]
et notons tout de suite que, comme $V$ est
de rang 2, on a un isom canonique
\[
\check V\simeq V\otimes(\det V)^{-1} \tag*{(52)}
\]
déduit de l'accouplement \deleted{symétrique} \inserted{alterné}
\[
\begin{aligned} V\otimes V&\longrightarrow\det V\\ x\otimes y&\longmapsto x\wedge y\end{aligned} \tag*{(53)}
\]
d'ailleurs la suite exacte (47) définit un
isom. canonique
\[
\det V\xrightarrow[\sim]{\alpha}L \tag*{(54)}
\]
\deleted{en associant à} par la condition
\[
\alpha(i_1(\lambda)\wedge x)=\lambda\pi_1(x)\qquad\lambda\in\Gamma(-,L),\quad x\in\Gamma(-,V) \tag*{(55)}
\]
En conjuguant (52) et (54), on trouve un
iso can
\[
\check V\xleftarrow[\sim]{\alpha}V\otimes L^{-1} \tag*{(56)}
\]
\deleted{par la condition} associé à l'accouplement \uuline{alterné}
\[
\left\{\begin{aligned}
(V\otimes L^{-1})\otimes V&\xrightarrow{\varphi}\underline O_S\\[-2pt]
\underset{(V\otimes V)\otimes L^{-1}}{\unsure{\Vert|\wr|\simeq}}&
\end{aligned}\right. \tag*{(57)}
\]
défini par la condition que
\[
\varphi(i_1(\underset{\substack{\rotatebox{-90}{$\in$}\\\Gamma L}}{\lambda})\otimes\underset{\substack{\rotatebox{-90}{$\in$}\\\Gamma V}}{x}\otimes\underset{\substack{\rotatebox{-90}{$\in$}\\\GammaV}}{\mu})=\lambda\mu\,\pi_1(x) \tag*{(58)}
\]
% archive sheet number: 35

\scanpage{37}
On peut encore dire que l'on a un
accouplement \uuline{alterné}
\[
 V\otimes V\xrightarrow{\varphi}\deleted{\unsure{K}}L \tag*{(57 bis)}
\]
défini par la condition que l'on a
\[
\varphi(i_1(\underset{\substack{\rotatebox{-90}{$\in$}\\\Gamma L}}{\lambda})\otimes\underset{\substack{\rotatebox{-90}{$\in$}\\\Gamma V}}{x})=\lambda\pi_1(x) \tag*{(58 bis)}
\]

Ceci dit, \deleted{\unsure{le fibré|la filt}} tensorisant (47) par $L^{-1}$
on trouve sur $L^{-1}\otimes V$ une structure
d'extension
\[
\unsure{\cdot|0}\longrightarrow\underline O_S\xrightarrow{i'}\underbrace{L^{-1}\otimes V}_{\substack{\wr\,\downarrow\,\alpha\\\check V}}\xrightarrow{\unsure{\pi|\pi'}}\deleted{\unsure{\mathcal O_S}}\check L\longrightarrow0 \tag*{(59)}
\]
% slanted
\margin{Généraliser cet énoncé au cas de la suite exacte (49) et de sa transposée\ldots}
qui, compte tenu de l'iso (56), n'est autre
que la structure d'extension (51) [i.e.
$\alpha$ transforme $i'_1$ en $i'$, $\pi'_1$ en $\pi'$].

Considérons l'hom composé
\diagram{\begin{tikzcd}[ampersand replacement=\&,column sep=small,row sep=small]
\& \mathrm{Sym}^n(\check V)\arrow[r,hook] \& \mathrm{Sym}^{*}(\check V)=\mathrm{Alg\,aff}\,\check{\mathbb V}(V)\arrow[r] \& \mathrm{Sym}^{*}(\check V)/(e_V-1)\arrow[d,equal]\\
\& L^{\otimes(-n)}\underline{\mathrm{Sym}}^n(V)\arrow[u,"\sim"',"\mathrm{Sym}^n(\alpha)=\alpha_n"]\arrow[d,"\sim"',"\mathrm{Sym}^n(\alpha)=\alpha_n"] \& \& \mathrm{Alg\,aff}(F)\\
\text{soit}\quad\text{(60)} \& \underline{\mathrm{Sym}}^n(\check V)\arrow[r,"j_n"]\arrow[d,equal] \& \mathrm{Alg\,aff}(F) \& \\
\& \substack{\text{faisceau des}\\\text{$n$-formes sur }V} \& \&
\end{tikzcd}}
% archive sheet number: 36

\scanpage{38}
\uline{Proposition} \uline{61} \inserted{$\forall n\in\mathbb N$} L'hom $j_n$ est injectif, et on
a commutativité dans le diagramme
\diagram{\begin{tikzcd}[ampersand replacement=\&,column sep=small,row sep=large]
L^{\otimes -n}\otimes\underline{\mathrm{Sym}}^{\unsure{n|*}}(V)\arrow[r,"\alpha_n","\sim"']\arrow[d,"i_n","{\mathrm{id}_{L^{\otimes(-n-1)}}\otimes\deleted{\unsure{i_n(L)}}}"'] \& \underline{\mathrm{Sym}}^n(\check V)\arrow[r,hook,"j_n"]\arrow[d,hook,"{\substack{\text{multiplication par}\\e_{\check V}=i'_1(1)}}"] \& \underline{\mathrm{Alg\,aff}}(F)\arrow[d,equal]\\
L^{\otimes(-n-1)}\otimes\mathrm{Sym}^{n+1}(V)\arrow[r,"\sim","\alpha_{n+1}"'] \& \underline{\mathrm{Sym}}^{n+1}(\check V)\arrow[r,hook,no head,"j_{n+1}"] \& \mathrm{Alg\,aff}(F).
\end{tikzcd}}
\margin{\deleted{\unsure{$i$}}}
Enfin, $\mathrm{Alg\,aff}(F)$ est isomorphe à la limite
inductive des $\mathrm{Sym}^n(\check V)\simeq L^{\otimes(-n)}\otimes\mathrm{Sym}^n(V)$.

\uline{NB} Les sections de $\mathrm{Sym}^n(\check V)$, identifiées à des
fonctions $\unsure{f|(-}$ sur $F$, s'appellent les fonctions \inserted{polynômiales} de
degré $\leq n$. Correspondance entre celles-ci et
formes sur \inserted{$V($\unsure{$=$}}$\check{\mathbb V}(V))$ de degré $n$. Pour \deleted{\unsure{tt|une}} telle
$f$, sa classe modulo $\mathrm{Sym}^{n-1}(\check V)$ (fonctions
polynômiales de degré $\leq n-1$) s'appelle
\uline{terme dominant} de $f$, et s'identifie
à une section de $L^{\otimes(-n)}$; \deleted{\unsure{c'est|ainsi}} \inserted{elle s'identifie} aussi à
la restriction à $L$ de la $n$-forme sur $V$ associée
à $f$. --- On notera que dans cette proposition
\margin{$!\,\left\{\vphantom{\begin{matrix}1\\1\\1\end{matrix}}\right.$}
se trouvent mélangés des \unsure{résultats} \inserted{valables} sur les
fibrés affines de rang quelconque, et ceux
spéciaux aux fibrés de rang 1, liés à
la dualité (52).
% archive sheet number: 37

\scanpage{39}
Soit maintenant (pour $n\geq1$ fixé) $D$ un
diviseur relatif de rang $n$ sur $F$. Soit
\inserted{$M_D\subset\underline{\mathrm{Sym}}^n(\check V)$} le sous-\deleted{\unsure{fibré de rang|faisceau des}}Module
des fonctions polynômiales sur $F$ de degré $\leq n$
\unsure{s'y|se|s'annulant} \unsure{tuant|trouvant|sur $D$}.

\uline{Proposition} $M_D$ est un supplémentaire de
$\underline{\mathrm{Sym}}^{n-1}(\check V)$ dans $\mathrm{Sym}^n(\check V)$, i.e. le composé
\diagram{\begin{tikzcd}[ampersand replacement=\&,column sep=small]
M_D\arrow[r,hook]\arrow[rrr,bend right=14] \& \mathrm{Sym}^{\deleted{\unsure{n-1}}n}(\check V)\arrow[r,dash] \& \mathrm{Sym}^n(\check V)/\mathrm{Sym}^{n-1}(\check V)\arrow[r,"\sim",no head] \& L^{\otimes -n}
\end{tikzcd}}
est un iso. L'application
\[
D\longrightarrow M_D
\]
est une bijection de l'ens des diviseurs
relatifs sur $F$ de degré $n$ avec l'ens
des sous-\deleted{\unsure{M}}Modules $M_D$ de $\mathrm{Sym}^n(V)$
satisfaisant la condition précédente. Cette
bijection est compatible avec les chgts de
base, et définit un isomorphisme
\diagram{\begin{tikzcd}[ampersand replacement=\&,column sep=small,row sep=small]
\underline{\mathrm{Div}}^{+n}_{F/S}\arrow[r,"\sim"] \& \substack{\text{\uline{\unsure{Splittages|Splittings} de}}\\
(0\longrightarrow\mathrm{Sym}^{n-1}(\check V)\longrightarrow\mathrm{Sym}^n(\check V)\longrightarrow L^{\otimes -n}\longrightarrow\unsure{0})}\\
\underset{\text{puissance sym. nième}}{\Sigma^n(F/S)}\arrow[u,"\sim"] \& \substack{\text{\uline{splittages de} } L^{\otimes\unsure{(+n)|n}}\otimes\text{cette suite}\\
\text{exacte, i.e. de la suite}\\ \inserted{\unsure{donnée|(50)}}\\
\deleted{\unsure{0\longrightarrow L\otimes\mathrm{Sym}^{n-1}(V)\longrightarrow\mathrm{Sym}^n(V)\longrightarrow\mathcal O_S\longrightarrow0}}}\arrow[u,no head,"\sim"]
\end{tikzcd}}
% archive sheet number: 38

\scanpage{40}
On trouve ainsi un isom canonique
\[
\begin{aligned}
\Sigma^n(F/S)\simeq{}&\text{fibré affine associé à l'extension}\\
&0\longrightarrow\mathrm{Sym}^{n-1}(V)\otimes L\longrightarrow\mathrm{Sym}^n(V)\longrightarrow\mathcal O_S\longrightarrow0\\
&\text{(dont le fibré des translations}\\
&\text{est }\mathrm{Sym}^{n-1}(V)\otimes L\text{ )}
\end{aligned}\tag*{(62)}
\]
Cette suite exacte s'interprète comme celle
\inserted{\unsure{associé|associée} à l'hom} allant des formes de degré $n$ sur $\check V$ \inserted{$\simeq V\otimes L^{-1}$} vers $\mathcal O_S$,
par la ``valeur en $e_{\check V}=e_B$''. Donc la donnée
d'une section de $\Sigma^n(F/S)$, i.e. d'un
diviseur relatif de degré $n$ sur $F$, équivaut
à la donnée d'une forme de degré $n$ sur
$B=\check V$, prenant la valeur 1 en $e_B$.

6. Spécialisons ceci au cas de $n=2$. On trouve

\uline{Théorème} Soit $F$ fibré affine, associé à
ext. (47) de $\mathcal O_S$ par ``\deleted{fibré}\inserted{\unsure{module|espace}} des translations'' sur $L$,
ou dualement par ext. $B$ de $\check L$ par $\underline O_S$, correspondant
à la donnée, dans le Module loc. libre
\inserted{$B$} de rang 2, \deleted{$B$} d'une section $e_B$ non nulle
en chaque fibre. Il y a alors une
corr. biunivoque entre

a) Les structures d'Algèbre comm. unifère
sur $B$ admettant $e_B$ comme \unsure{section|Slection|élection|élément} unité.

b) Les diviseurs relatifs \inserted{\deleted{\unsure{$C$}} \unsure{$\Omega$|de} sur $F$} de degré 2 sur $S$

c) Les sections de $\Sigma^2(F/S)$
% archive sheet number: 39

\scanpage{41}
\margin{\rule{0.35pt}{2.2\baselineskip}}d) Les formes quadratiques $q$ sur $B$ satisfaisant
\deleted{\unsure{les}}la condition
\[
 q(e_B)=1 \tag*{(63)}
\]

Les équivalences entre b), c), d) viennent
d'être explicitées, celle entre a) et b) a
été vu dans 5). [Il reste encore à expliciter
les relations directes entre a) et d), et rappeler
(à propos de l'équiv. a) et b)) la relation
\[
\Omega=\mathrm{Spec}(B) \tag*{(64)}
\]
lorsque $\Omega$ \inserted{$\subset F$} et une structure d'Algèbre sur $B$
se correspondent. Enfin, on aura
\[
q=\unsure{\mathcal N|\mathscr N|N}_{B/\mathcal O_S} \tag*{(65)}
\]

[[Pour vérifier ce dernier point, nous allons le
reformuler pour $n$ quelconque. Notons que
la donnée d'une section de $\Sigma^n(F/S)$, ou
encore d'une $n$-forme $f_n$ sur $B$ telle que
$f_n(e_B)=1$, ou encore d'un $n$-diviseur relatif
effectif $\Omega_n$ sur $S$, équivaut à la donnée d'une
$\mathcal O_S$-Algèbre \inserted{\unsure{commu.|comm.}} $C_n$ \inserted{loc. libre} de rang $n$, et d'un hom \inserted{surj}
\[
\underset{\substack{\wr\downarrow\\\underline{\mathrm{Sym}}^*(B)/(e_B-1)}}{\underline{\mathrm{Alg\,Aff}}(F)}\longrightarrow C_n
\]
ou ce qui revient au $=$, d'un hom de \inserted{$\mathcal O_S$-Mod.}
\[
 B\xrightarrow{g} C_n \tag*{(66)}
\]
% archive sheet number: 40

\scanpage{42}
\unsure{tel|tels} que
\[
\left\{\begin{aligned}\text{a) }&g(e_B)=1_{C_n}\\
\text{b) }&g(B)\text{ engendre }C_n.
\end{aligned}\right. \tag*{(67)}
\]
Ceci dit, considérons \deleted{l'application \unsure{fonction}} \inserted{le morphisme}
\[
F_n:\check{\mathbb V}(C_n)\longrightarrow\check{\mathbb V}(\mathcal O_S) \tag*{(68)}
\]
défini par
\[
 F_n(x)=\unsure{\mathcal N|\mathscr N|N}_{C_n/\mathcal O_S}(x)\;(=\mathcal N_{\Omega/S}(x)), \tag*{(69)}
\]
c'est une $n$-forme, et on peut considérer
sa ``restriction'' $F_n\circ g$. Ceci dit, on
aura (sauf erreur)
\[
 f_n=\deleted{\unsure{(-1)^n}}F_n\circ g\ \inserted{=\deleted{\unsure{(-1)^n}}\mathcal N_{C_n/\mathcal O_S}\circ g}\quad\text{dans }\Gamma\underline{\mathrm{Sym}}^n(\check B=V) \tag*{(70)}
\]
[ce qui montre que la correspondance
entre $n$-diviseurs sur $F$, \deleted{et} i.e. syst. $(C_n,g)$
satisfaisant (67), et $n$-formes définie par
(70), est bijective].

Pour vérifier (70), OPS que $B$ est muni d'
une section $t$ telle que $(e_B,\unsure{T|t})$ forment une
base. Alors $C_n$ sera engendré sur $\mathcal O_S$ par \inserted{l'image $t$ de} $T$,
donc aura la base $1,t,\ldots,t^{n-1}$, et sera défini
par
\[
 C_n\simeq\mathcal O_S[T]/(f)\mathcal O_S[T] \tag*{(71)}
\]
où $f$ est un polynôme unitaire \inserted{en $T$} de \uline{degré} $n$,
\[
 f(T)=\deleted{\unsure{a_0+}}T^n+a_1T^{n-1}+a_2T^{n-2}\,\text{---}\,+a_n
\]
\deleted{qui n'est autre que $f_n(t)$}
On \deleted{\unsure{a}} vérifie de suite que l'on a $N$
\[
\begin{aligned}
 M_n&=\mathcal O_S\,g_n\qquad\deleted{\unsure{g_n=e_B^n+a_1 e_B^{n-1}T+\cdots+a_n\in\Gamma\mathrm{Sym}^n B|q_n=e_B^n+a_1e_B^{n-1}T+\cdots+a_n\in\Gamma\mathrm{Sym}^n B}}\\
 \deleted{(72)}\quad&\deleted{\unsure{f_n(u e_B+vT)=u^n+a_1u^{n-1}v+a_2u^{n-2}v^2+\cdots+a_nv^n}}\\
 &\unsure{B}\quad g_n=\deleted{\unsure{a_n}}T^n+a_{n-1}e_BT^{n-1}+\cdots+a_ne_B^n\in\Gamma\mathrm{Sym}^n B
\end{aligned}\tag*{(72)}
\]
% archive sheet number: 41

\scanpage{43}
d'où on conclut \unsure{(73)} $\nearrow$
et la formule à vérifier devient
\[
\begin{aligned}
\unsure{\mathcal N|\mathscr N|N}_{C_n/\mathcal O_S}(u1_C+vt)&=\deleted{\unsure{(-1)^n}}u^n-a_1u^{n-1}v\,\text{---}\,+\inserted{(-1)^n}a_nv^n\\
&=(-1)^nv^nf\left(\deleted{\unsure{+}}-\frac{u}{v}\right)
\end{aligned}\tag*{(\deleted{\unsure{73}}74)}
\]
Or si $u=0$, $v=1$ cette formule se réduit à
\[
 \mathcal N_{C_n/\mathcal O_S}(t)=\inserted{(-1)^n}a_n \tag*{(\deleted{\unsure{74}}75)}
\]
Or la matrice de la multiplication par $t$, dans
la base $1,t,\ldots,t^{n-1}$, \unsure{est}
\[
\begin{pmatrix}
0&0&&0&-a_n\\
1&0&&0&-a_{n-1}\\
0&1&&0&\\
\vdots&\vdots&&\vdots&\vdots\\
0&0&\cdots&0&\\
0&0&&1&-a_{\unsure{1|4}}
\end{pmatrix}
\]
dont le déterminant est $\underbrace{(-1)^{n+1}}_{\substack{\text{signature de la}\\\text{\deleted{\unsure{permuta}}permutation circulaire}\\\text{de l'ens à $n$ élts}}}(-a_n)=(-1)^{n\deleted{\unsure{+1}}}a_n$.

\deleted{Il suffit \unsure{pour} cela d'un} \inserted{On se ramène au cas de $t$ (au lieu de $u+vt$)}
par un argument universel. On peut aussi le
prouver directement, sans passer par des
déterminants, en faisant une ext. fid. plat
de la base permettant d'écrire
\[
 f\deleted{\unsure{_n}}(T)=\prod_{i=1}^n(T-t_i)\qquad\text{donc }a_i=(-1)^i\unsure{\tau|\sigma}_i(t_1,\ldots,t_n) \tag*{(76)}
\]
avec des sections $t_i$ de $\underline O_S$ $(1\leq i\leq n)$. \deleted{\unsure{On a}}
Alors (75) devient
\[
\begin{aligned}
\mathcal N_{C_n/\mathcal O_S}(u1_C+vt)&=(-v)^nf\left(\deleted{1}-\frac{u}{v}\right)\deleted{\unsure{v}}\;=\prod(u\deleted{\unsure{T}}+vt_i)\deleted{\unsure{(-1)^n}}\\
&=\prod_i \rho_i(u1_C+vt)
\end{aligned}\tag*{\deleted{\unsure{(77)}}}
\]

\noindent\rule{\linewidth}{0.3pt}

utilisant $\check V\xleftarrow{\sim} V\otimes L\simeq V$ \qquad $e_B\longmapsto T'$, $T\longrightarrow -e_B'$
(ici $(e_B',T'$ est la base duale de $e_B,T$)
% slanted
\margin{facteur de proportionnalité}
\[
 f_n=\underbrace{\varepsilon}_{\leftarrow}\deleted{\unsure{\lambda}}(-1)^n\left({e_B'}^n-a_1{e_B'}^{n-1}T'+\cdots+(-1)^ia_i{e_B'}^{n-i}{T'}^i\cdots+(-1)^na_n{T'}^n\right) \tag*{\deleted{\unsure{(73)}}}
\]
i.e.
\[
 f_n(ue_B+vT)=\deleted{\unsure{\varepsilon(-1)^n|\varepsilon}}\left(u^n-a_1u^{n-1}v+\cdots+(-1)^ia_iu^{n-i}v^i\,\text{---}\,+(-1)^na_nv^n\right) \tag*{(73)}
\]
(avec $\varepsilon=(-1)^n$ pour avoir $f_n(e_B)=1$)
% archive sheet number: 42

\scanpage{44}
où
\[
\rho_i:C_n\longrightarrow\underline O_S \tag*{(77)}
\]
est l'hom de $\underline O_S$-Algèbres défini par
\[
\rho_i(T)=t_i, \tag*{(78)}
\]
et alors la formule à prouver est \unsure{cas|un|un cas} particulier
de la suivante, bien connue
\[
 \unsure{\mathcal N|\mathscr N|N}_{C(t_1,\ldots,t_n)/\mathcal O_S}(\deleted{\unsure{f}}w)=\prod_{1\leq i\leq n}\rho_i(w)\quad ]] \tag*{(79)}
\]

Revenons \deleted{\unsure{au théo|au théorème}} au cas $n=2$, et à la correspondance
entre structures d'Algèbres quadratiques
sur $B$, admettant $e_B$ comme él\textsuperscript{t} unité, et
\deleted{structures \inserted{de} \unsure{quad|quadr}} formes quadratiques $q$ sur $B$ satisfaisant
$q(e_B)=1$. Explicitons la structure sur
$\Sigma^2(F)$ de fibré \deleted{vectoriel}\inserted{affine} sur $2F$, de \deleted{base}
module des translations $(L^{\otimes2})_{2F}$. Il
faut d'abord expliciter
\[
\underline b(q)\in\Gamma(2F)\subset\Gamma(\check B=V) \tag*{(80)}
\]
pour $q$ comme dessus. \deleted{\unsure{On fera|On fixe}}

\uline{Prop.} Soit $\varphi$ la forme bilinéaire associée à
$q$, et considérons la forme
\[
\left|\begin{aligned}
\underline b(q)&:B\longrightarrow\underline O_S\\
\underline b(q)(x)&=\varphi(e_B,x)\deleted{\unsure{+}}\;(=\mathrm{tr}_{B/\mathcal O_S}\,x)
\end{aligned}\right. \tag*{(81)}
\]
On a $\underline b(q)(e_B)=\varphi(e_B,e_B)=2q(e_B)=2$, donc
$\underline b(q)$ est une section de $2F\subset\check{\mathbb V}(\check B)$, qui n'est
% archive sheet number: 43

\scanpage{45}
autre que $\underline b(\Omega)$ du \S~5.

Fixons un $\underline b(q)$ \inserted{$=\mathrm{tr}$} i.e. \deleted{\unsure{$q$}} \uline{la forme} $\mathrm{tr}$
\uline{trace} $B\longrightarrow\mathcal O_S$, et étudions les $q$ correspondants
à la forme trace donnée. On trouve
que ceux-ci \deleted{\unsure{font un|faisceau}} \inserted{sont les sections d'un faisceau qui est un} \deleted{\unsure{pré}}torseur (et $=$ un
torseur\ldots) \unsure{sous|sur|à} le faisceau ou Module des
formes quadratiques $q_0$ telles que
\[
\left\{\begin{aligned}q_0(e_B)&=0\\
\varphi_0(e_B,\unsure{-|\cdot})&=0\qquad\deleted{\forall x}
\end{aligned}\right.\quad
\substack{\text{(\unsure{où|ici} $\varphi_0$ est la}\\\text{forme bil associée}\\\text{à $q$ )}} \tag*{(82)}
\]
\deleted{l'hom $V\longrightarrow\check V$}
\deleted{i.e. $\varphi_0(e_B,\unsure{-})$ une forme $\varphi_0$}
ce sont aussi les formes quadratiques
qui se factorisent par $B\longrightarrow\check L$, i.e.
qui correspondent aux formes quadratiques
sur $L$, ou encore aux \deleted{\unsure{fo|for}} sections de
$L^{\otimes2}$. Donc les formes en question
forment \deleted{les} sections\deleted{\unsure{en}} d'un fibré affine sous
$L^{\otimes2}$, dont les sections $c_0$ correspondent aux $q_0$ par
\[
q_0(x)=c_0\,\pi_1(x)^2. \tag*{(83)}
\]
% archive sheet number: 44

\scanpage{46}
7. \deleted{\unsure{Symétrie|Symétries}} Involution canonique.

\uline{Prop} Soit $B$ Algèbre quadratique sur $\mathcal O_S$.
\margin{\rule{0.35pt}{3\baselineskip}}Alors l'hom $\mathcal O_S$-linéaire
\[
\sigma:u\longmapsto\deleted{\unsure{\mathrm{tr}\,u\,1_B}}=(\mathrm{tr}\,u)e_B-u \tag*{(84)}
\]
est un \uline{automorphisme} d'algèbres involutif,
caractérisé par la condition
\[
u+\sigma u=\mathrm{tr}\,u\,e_B \tag*{(85)}
\]
% margin note written sideways
\margin{On \unsure{appellera|appelle|veut} quasi-réflexion d'un module $B$, rel. à un sous-module $B'$ tel que $B/B'\simeq\check L$ inversible, tt endomorphisme \deleted{\unsure{par}} de la forme $\mathrm{id}+v$, où $v\in L\otimes V\subset\check V\otimes V$ est tel que $\mathrm{tr}\,v=-2$.}

On notera que la donnée (pour $B$ un $\mathcal O_S$-Module
fixé avec $e_B$ \inserted{section} \deleted{\unsure{partout}} non nulle en chaque fibre)
la donnée de l'involution $\sigma$ équivaut à celle
de la forme \inserted{$\underline b=$} $\mathrm{tr}$\deleted{\unsure{$b$}} (section de $2F$), \deleted{\unsure{et nous}}
Moyennant cette correspondance, les $\underline b$
correspondent \deleted{aux} exact\textsuperscript{t} aux \inserted{quasi-}\uline{réflexions} $\sigma$
du $\mathcal O_S$-Module $B$ telles que $\sigma e_B=e_B$, i.e.
admettant $\mathcal O_S\cdot e_B$ comme droite fixe\unsure{/} \deleted{\unsure{soit}}

% slanted
\margin{\deleted{\unsure{Sauf erreur|Sauf en faisant|déf en faisant} si \unsure{on est partout|au point} \unsure{de car. résiduelle|car si résiduellement} \unsure{char|de} 2, $\mathrm{tr}\,\unsure{v|u}$ n'est pas \unsure{forcément nulle|partout nulle} sur les fibres résiduelles\ldots}}
\deleted{Cette réflexion \unsure{n'est effet|n'est|n'est effect.}, un endomorphisme
ayant $\mathcal O_S e_B$ comme droite de pts fixes se met
sous la forme générale}
\[
\deleted{\sigma u=u+a\deleted{\unsure{\otimes}}\pi_1(u)\qquad\text{i.e. }\sigma=\mathrm{id}_B+\deleted{\unsure{a\otimes}}\pi_1\otimes a} \tag*{\deleted{\unsure{( }}}
\]
\deleted{où $a\in\Gamma(B\otimes L\simeq V=\check B)$ et \deleted{\unsure{où}} (et $\pi_1\in\Gamma(\check B\otimes\check L\simeq\underline{\mathrm{Hom}}(B,\check L)\subset B)$ \inserted{pour mém\textsuperscript{re}}}
\[
\begin{aligned}
\deleted{\sigma^2}&\deleted{=\mathrm{id}_V+2\pi_1\otimes a+\deleted{\pi_1(a)}\inserted{\langle a,\pi_1\rangle}\pi_1\otimes a}\\
&\deleted{=\mathrm{id}_V+\pi_1\otimes a(2+\langle a,\pi_1\rangle)}
\end{aligned}
\]
\deleted{donc $\sigma$ est involutive \deleted{\unsure{ssi}} et \deleted{\unsure{part|partout}} $\ne\mathrm{id}$ en tte
fibre i.e. est une réflexion propre, ssi $a$ \deleted{est} $\ne0$
en tte fibre, \uline{et} $\langle a,\pi_1\rangle\deleted{\unsure{a}}=-2$. \deleted{\unsure{voir}} Identifions}
% archive sheet number: 45

\scanpage{47}
\deleted{$a$ à une forme linéaire sur $B$, et $\pi_1$ à la
section $e_B$ de $B$ (au signe près\ldots) cette relation
s'écrit $a(e_B)=-2$, et \deleted{cette équation}
on vérifie (\deleted{\unsure{sachant bien}} \inserted{faisant attention} au \unsure{traquenard|traitement} des
signes) que $a$ correspond à $\underline b=\mathrm{tr}$ \ldots\ Le
seul cas est celui où}

Cette quasi-réflexion est une
réflexion i.e. est $\ne$ identité en chaque fibre, \unsure{sss|ssi}
en tt pt $s\in S$ de car résiduelle 2, on a
$\mathrm{tr}_s$ sur $B_s$ non (id\textsuperscript{t}) nulle. Ceci résulte \unsure{des|de} la

\uline{Proposition} Pour qu'une quasi-réflexion d'un
module \deleted{\unsure{quadratique, soit rg 2}} \inserted{\unsure{quelconque|partout de}} \deleted{de} $B$, (par r. à un
$B'$ tel que $B/B'\simeq\check L$ inv. \uline{donné}) soit l'identité,
il f et s. que $S$ soit de car. 2 (i.e. $2.1_S=0$)
et que la \deleted{\unsure{$f$}} section de $L\otimes B$ qu'elle définit soit
nulle \ldots\ donc dans le cas actuel ($B$ de rg 2,
$B'=\mathcal O_S e_B$, donc $L\otimes B\simeq\check B$, que la forme $\underline b=\mathrm{tr}$
sur $B$ qu'elle définit (et qui la définit, car
$\underline b(e_B)=2$ \unsure{---|\ldots}) soit nulle.

\deleted{Supposons}\inserted{Revenons} maintenant au cas où $\underline b=\mathrm{tr}$ \deleted{\unsure{provient}} \inserted{est la}
\inserted{forme trace} d'une structure d'Algèbre sur $B$ ayant $e_B$ comme
unité. On obtient alors la formule bien connue
\[
u.\sigma u=\unsure{\mathcal N|\mathscr N|N}(u)e_B=q(u)e_B \tag*{(86)}
\]
i.e. \qquad $u(\mathrm{tr}\,u-u)=\mathcal N(u)$ \qquad ou encore
\[
\big|\,u^2-u\,\mathrm{tr}\,u+\mathcal N(u)=0 \tag*{(87)}
\]
% archive sheet number: 46

\scanpage{48}
qu'on peut interpréter comme l'équation de
Hamilton-Cayley pour la \inserted{l'op. de} translation \deleted{\unsure{induite}} définie
par $u$ dans $B$.

\uline{8} \uline{Cas où 2 inversible sur} $S$ (car résiduelle $\ne2$)

Dans ce cas la donnée d'une section $\underline b=\mathrm{tr}$ de
$2F$ i.e. d'une forme $\mathrm{tr}$ sur $B$ telle que $\mathrm{tr}(e_B)=2$,
équivaut, \deleted{\unsure{(en en divisant par} \inserted{en prenant} (en posant $\tfrac12\underline b=\tfrac12\mathrm{tr}\,u\unsure{:|=}\overset{\unsure{s.}}{\unsure{s}}$}
\[
 s_0=\tfrac12\deleted{\unsure{\mathrm{tr}|\underline b}}\underline b=\tfrac12\mathrm{tr}\,\unsure{u|\ldots} \tag*{(88)}
\]
à la donnée d'une section \inserted{origine} $s_0$ de $F$, ou encore
d'un splittage de l'extension $B$ de $\check L$ par $\mathcal O_S$ en
un iso
\[
\alpha_{s_0}:\underline O_S+\check L\xrightarrow{\sim}B \tag*{(89)}
\]
Le sous-module $\alpha_{s_0}(\check L)$ de $B$, supplémentaire
de $\mathcal O_S$, n'est autre que le \uline{noyau} de la
forme $\mathrm{tr}$ (qui est aussi celui de la forme
$\tfrac12\mathrm{tr}$ \ldots), \inserted{ou encore le sous-espace propre de $\sigma$ pour val. propre $-1$.}
% slanted
\margin{NB en termes de $q=\unsure{\mathcal N|\mathscr N|N}_{B/\mathcal O_S}$, $\alpha_{s_0}(\check L)$ est l'\uline{orthogonal} de $e_B$, i.e. de $\mathcal O_S e_B$.}
Ceci posé, la donnée d'une
structure d'Algèbre sur $B$ compatible avec
$(e_B,\mathrm{tr})$ équivaut à la donnée de
$\underline b(s_0)$\inserted{$\in\Gamma L$}, $\underline c(s_0)$\inserted{$\in\Gamma L^{\otimes2}$} correspondant au scindage $s_0$.
Or \inserted{\unsure{(14))|(14)}} \deleted{\unsure{au}} $\underline b_{\unsure{0}}=\underline b(s_0)$ est la forme trace restreint à
$\alpha_{s_0}(\check L)$, et dire que \deleted{$\underline b(s_0)$} $s_0$ est le scindage canonique
signifie exactement que
\[
\deleted{\unsure{b=}}\underline b(s_0)=0,\tag*{(90)}
\]
% archive sheet number: 47

\scanpage{49}
de sorte que l'extension quadratique est décrite
entièrement par la section $c_0$ de $L^{\otimes2}$. Or comme
\[
\deleted{\unsure{c_0=\mathcal D}}\ \unsure{\mathcal D|\mathscr D|\mathfrak d|\partial}\deleted{\unsure{(|{}_{\xi}}}=b_0^2-4c_0=-4c_0\tag*{\deleted{\unsure{a|$\underline c$}}(91)}
\]
\deleted{\unsure{on a aussi|et aussi}}
la donnée de $c_0$ équivaut à celle de $\mathcal D$. On
trouve ainsi

\uline{Th} Si $2\,1_S$ est inv., la catégorie des Algèbres
quadratiques sur $\mathcal O_S$ (ou \deleted{celle} des rev. quadratiques
\inserted{$\Omega$}\deleted{\unsure{de}} de $S$) équivaut à celle des couples $(L,\mathcal D)$
\uline{où} $L$ est un $\mathcal O_S$-Module inversible, et $\mathcal D$
une section de $L^{\otimes2}$. \deleted{\unsure{Le|Lorsque}} Le fibré affine \deleted{\unsure{de rang 1}}\inserted{enveloppant $\Omega$}
\deleted{\unsure{$E$|$\xi$}} est can\textsuperscript{t} scindé, et son Module des translations
est $L$, de sorte que $\Omega\subset\check{\mathbb V}(L)$\deleted{\unsure{)|$\xi$}}. Ceci
dit, si \deleted{\unsure{$\Omega$ s'identifie|l'Algèbre}}\inserted{$\Omega(\mathcal D)$ est le} \deleted{\unsure{le}} sous-schéma de $\check{\mathbb V}(L)$
des sections $x$ satisfaisant $x^2=\mathcal D$, on a \deleted{\unsure{$\Omega=\frac12$}}
\[
\Omega=\frac12\Omega(\mathcal D)\tag*{(92)}
\]

Il reste à vérifier cette dernière assertion.
Or pour un scindage $s$ quelconque, si $\underline b,\underline c$ sont
les sections associées de $L,L^{\otimes2}$, On a vu
dans \S\,2 (variante marginale) que $\Omega$ s'identifie
au sous-schéma de $\check{\mathbb V}(L)$ d'équation
\[
x^2-bx+c=0\tag*{\unsure{(6)|(5)|(*)}}
\]
Ce qui ici, \deleted{\unsure{si}} comme $b=b_0=0$, s'écrit
\[
x^2=\frac14\mathcal D
\]
ou encore
\deleted{\illegible{1}}
\[
(2x)^2=\mathcal D
\]
% archive sheet number: 48
